\documentclass[12pt]{article}
\usepackage[T1]{fontenc}
\usepackage[utf8]{inputenc}
\usepackage{lmodern}
\usepackage{textcomp}
\usepackage{enumitem}
\usepackage{geometry}
\usepackage{graphicx}
\usepackage{amsmath, amssymb}
\geometry{margin=1in}
\begin{document}
\begin{center}
Psychology - Graduate Level Assessment 7 \
Total Marks: 40 \
Answer all questions.
\end{center}

\section*{Multiple Choice (10 marks)}
\begin{enumerate}[label=\arabic*.]
\item Exposing individuals to arguments against their views, arguments that are then strongly refuted, may serve to
\begin{enumerate}[label=(\alph*)]
    \item make them more argumentative
    \item make them more susceptible to accepting opposing views
    \item weaken their resistance to later persuasive appeals
    \item increase their attitudinal ambivalence
    \item strengthen their existing views
\end{enumerate}
\item According to Harry Stack Sullivan, which of the following cognitive experiences characterizes schizophrenic behavior:
\begin{enumerate}[label=(\alph*)]
    \item Parataxic Distortion Mode: dealing w/ other like they're from early life
    \item Parataxic Integration Mode: connecting unrelated events or ideas
    \item Basic anxiety: feeling of helplessness/isolation in a hostile world
    \item Eudaimonic Mode: pursuit of meaning and self-realization
    \item Autistic Mode: living in a world of fantasies and daydreams
\end{enumerate}
\item According toJellinek(1952) what progression does thelife ofan alcoholic usually follow ?
\begin{enumerate}[label=(\alph*)]
    \item Onset phase, abuse phase, dependency phase, abstinence phase
    \item Prealcoholic phase, alcoholic phase, recovery phase, relapse phase
    \item Initial phase, adaptive phase, dependent phase, deteriorative phase
    \item Experimental phase, social use phase, intensive use phase, compulsive use phase
    \item Initial phase, middle phase, crucial phase, termination phase
\end{enumerate}
\item Discuss the two main theoretical points of view on the origins of human aggressive behavior.
\begin{enumerate}[label=(\alph*)]
    \item Biological predisposition and environmental triggers
    \item Evolutionary psychology and learned behavior theory
    \item Classical conditioning and operant aggression
    \item social learning theory and innate aggression
    \item Hormonal imbalance theory and cultural aggression
\end{enumerate}
\item With regard to minority and nonminority clients, psychotherapy is
\begin{enumerate}[label=(\alph*)]
    \item more effective for minority clients
    \item only effective for nonminority clients
    \item more effective when client and therapist have different racial/ethnic origins
    \item equally effective
    \item only effective when client and therapist speak the same language
\end{enumerate}
\end{enumerate}

\section*{True / False (10 marks)}
\textit{Answer True or False}
\begin{enumerate}[label=\arabic*.]
\setcounter{enumi}{5}
\item True or False: The mean of the given data set is 845.45, the median is 900, and the mode is 900.
\item True or False: Education is always impartial, while propaganda consistently conveys biased information to its audience.
\item True or False: Electrical stimulation can effectively close the gate in the gate-control theory of pain perception.
\item True or False: Job productivity is a significant predictor of employee retention over long periods in the workplace.
\item True or False: Participation in research is ethical as long as students are not coerced into taking part in the study.
\end{enumerate}

\section*{Long Form Response (20 marks)}
\textit{Answer each question with a clear and complete explanation.}
\begin{enumerate}[label=\arabic*.]
\setcounter{enumi}{10}
\item Discuss the psychological implications and risks associated with withdrawal delirium, particularly focusing on the substance most commonly linked to severe withdrawal symptoms.
\item Discuss the cultural values associated with Asians and Asian Americans, and critically analyze why fatalism is considered least characteristic within this context.
\end{enumerate}

\vfill
\noindent\textit{End of Paper}
\end{document}